% !TeX root = ../main.tex
\section{Limitations}
\label{sec:limitations}

\paragraph{The generations differ in several ways at once.}
This is the central limitation and it bounds every causal claim in the chapter. Between
generation~1 and generation~3 the therapist changed (Llama-2-7B $\to$ Llama-3.2-1B), the
patient simulator changed (GPT-3.5 $\to$ \primaryjudge{}), the patient personas became less
cooperative, the rubric moved from regex parsing to a JSON schema, a minimum-context filter was
added, and the preference-tree construction was revised. The chapter establishes that the
look-ahead effect reverses across these settings and rules out two specific explanations
(analysis method, \S\ref{sec:gen1}; grader identity, \S\ref{sec:regrade}). It \textbf{cannot}
attribute the reversal to model capacity, patient difficulty, or oracle strength individually.
The moderator of \S\ref{sec:discussion} is a reading that survives the evidence, not a
demonstrated cause, and it rests on $n = 5$ arms with one contradicting case.

\paragraph{Every result here is within PTO.}
All three generations optimize with PTO, which is what makes them a comparable series, but it
also means ``look-ahead never leads'' is established for preference-tree optimization and not as
a statement about look-ahead in general. This is a real restriction, because $\K$ is
optimizer-agnostic by construction --- it changes what the oracle sees, not how candidates are
weighted (\S\ref{sec:pto}) --- so the mechanism of \S\ref{sec:mechanism} predicts the same
outcome under a group-relative loss, and that prediction is untested. Running a matched $\K$
comparison under a second optimizer is the single highest-value addition to this chapter.

\paragraph{Only one non-zero look-ahead depth.}
All three generations compare $\K{=}0$ against $\K{=}5$. Generation~1 also ran a $\K{=}3$ sweep,
but with different hyperparameters (a different filter threshold and sampling temperature), so
it is not matched and is excluded here. Nothing in this chapter speaks to whether a different
depth --- shorter, longer, or adaptive --- would behave differently, and the mechanism of
\S\ref{sec:spread} suggests the failure is about what the extra turns contain rather than how
many there are.

\paragraph{Absolute scores are not comparable across generations.}
Generation~2 generated conversations in 4-bit NF4 and generation~3 in bf16, which changes
degeneration rates by roughly $30\times$ and depresses generation~2's absolute levels by about
$0.6$ points for the same base model. Every claim here is therefore a \emph{within-generation}
paired contrast. Cross-generation comparisons of level --- including the vertical positions in
Figure~\ref{fig:trajectories} --- are not interpretable, and the moderator analysis uses gains
and slopes relative to each generation's own base for this reason.

\paragraph{The evaluation is entirely simulated.}
Patients are language models, and the referee is a language model filling in questionnaires
designed for humans. Generation~3 buys down part of this risk with a decoupled second grader
that reproduces the $\K$ result, and the oracle's own reproducibility is measured
($\mathrm{ICC}(2,1) = 0.86$--$0.99$; $0.96$--$0.99$ on Q1 and Q2). But no human rated any
conversation in any generation, and nothing here establishes that a higher questionnaire score
corresponds to better therapy for a real client. The chapter's claims are about a lever's effect
on a measured quantity, not about clinical benefit.

\paragraph{Generation 1 was re-graded on Q1 and Q2 only.}
The re-scoring pass (\S\ref{sec:one-grader}) covered the two questionnaires all three generations
share. Generation~1's conversations were not scored on the wider battery (MITI, change-talk,
MI-inconsistency), so the cross-generation comparison is restricted to the satisfaction and
working-alliance axis. Generation~2's $105$-contrast sweep does cover six questionnaires, and
generation~3's covers eight; neither shows a look-ahead effect on any of them.

\paragraph{One re-grading draw.}
Generation~1 was re-scored once per conversation rather than with repeated draws. Given measured
oracle reproducibility on Q1/Q2 and $96$ paired personas per cell, sampling noise on an arm mean
is roughly $0.01$ --- small against the effects discussed --- but the re-grade carries no
per-cell reliability estimate of its own.

\paragraph{Statistical multiplicity.}
Holm correction is applied within each (generation, oracle/optimizer, iteration) family across
metrics, not globally across the roughly $150$ contrasts in the chapter. This is the correction
used by generation~3's existing artifacts and is kept for consistency. It is mildly liberal at
the chapter level; note that the direction of that liberality favours \emph{finding} effects,
and the chapter's principal claims are nulls and reversals, which such a correction makes harder
rather than easier to assert.

\section{Conclusion}
\label{sec:conclusion}

Look-ahead simulation --- branching a conversation, playing it forward under the current policy,
and scoring where the reply leads rather than the reply itself --- was the component our ICLR
2025 paper credited for its best results. Across three generations of the same experiment,
re-analysed with one set of statistics and placed on one measurement axis by re-grading the
oldest generation with the newest oracle, that benefit does not survive. It is large and
consistent in generation~1, absent across $105$ paired contrasts in generation~2, and never
present across eight matched iterations and two independent graders in generation~3.

The original result was not a statistical artifact --- a stricter test strengthens it --- and it
was not a grader artifact, since the modern oracle reproduces it on the same conversations. What
changed was the setting, and the property of the setting that best organizes all three
generations is how well the myopic reward works without help. Look-ahead in this domain buys a
rescaling of the reward rather than more information about it: no extra branch points, an
unchanged margin-to-noise ratio, no added proxy faithfulness at matched policy, and no benefit
from the $70\%$ of extra preference pairs it generates. Where the one-step signal is inert, a
rescaling is enough to get a run moving. Where it already works, look-ahead is several times the
compute for nothing --- and, at the end of the longest run we have, slightly worse than nothing.
